\section[Proceso gauge euclídeo Yang--Mills en dimensión 4]{Construcción
proyectiva de un proceso gauge euclídeo Yang--Mills, reconstrucción de
Osterwalder--Schrader y brecha colectiva en dimensión \(4\)}


\noindent La especialización de la dinámica nonádica enriquecida,
construida en las secciones precedentes, es
\[
 (\widehat X_\infty,\Gamma_9)
 \xrightarrow{\ \mathcal R_{\rm YM}\ }
 (\mu_m,T_m,H_m,\Omega_m)
 \longrightarrow(\mu_\infty,\mathrm{OS},E(4))
 \longrightarrow\operatorname{gap}(H_{\rm phys})>0.
\]
El lector clover que expresa \(F^2\) reconoce la curvatura sobre este mismo
proceso proyectivo; no construye retrospectivamente \(\mu_m\). La teoría
gauge, la reconstrucción OS, el vacío simple y la brecha ya quedan fijados
por esta cadena. La identidad
\eqref{eq:amp-ym-identidad-accion-requerida} compara después una
parametrización opcional por \(g_R\); es un control no gobernante y no una
condición para identificar la teoría Yang--Mills construida.

\subsection*{Objeto gauge y dominio de refinamiento}

Sea \(G\) un grupo de Lie compacto, conectado y simple. Para cada
\(m\geq0\), se considera la red euclídea
\begin{equation}\label{eq:realizaciones-ym-red}
 \Lambda_m=a_m\mathbb Z^4,
 \qquad a_{m+1}=\frac{a_m}{9}.
\end{equation}
La álgebra cilíndrica gauge se genera por variables de enlace en \(G\),
operadores de entrelazamiento de representación y las coordenadas
enriquecidas de color,
orientación, holonomía, acarreo de fusión, ruta y memoria. Sea \(\mu_m\) la
medida invariante común y sea
\[
 \mathcal H_m=L^2(\mathcal A_m,\mu_m)^{\rm gauge}.
\]
El vector constante \(\Omega_m\) define
\begin{equation}\label{eq:realizaciones-ym-proyectores}
 P_{\Omega_m}=|\Omega_m\rangle\langle\Omega_m|,
 \qquad Q_m^{\rm phys}=I-P_{\Omega_m}.
\end{equation}
El segundo proyector es el complemento físico del vacío; no se identifica
con un residual empleado únicamente para organizar la escala de
renormalización.

El circuito de actualización es local:
\begin{equation}\label{eq:realizaciones-ym-circuito}
 K_m^{\rm loc}
 =\prod_{r=1}^{9\cdot81}
   \prod_{B\in\mathcal B_{m,r}}K_{m,B}.
\end{equation}
La partición \(\{\mathcal B_{m,r}\}_r\) enumera cada coordenada una sola vez; los factores
del circuito se identificarán abajo con el semigrupo de sus expectativas
condicionales, y no con nueve dinámicas independientes.
Los bloques de un mismo color poseen soportes disjuntos y el diámetro físico
de cada soporte está uniformemente acotado al variar \(m\). Esta localidad
forma parte del objeto, de modo que el refinamiento no convierte una
actualización elemental en una operación instantáneamente global.

Sea \(T_m\) la transferencia inducida por
\eqref{eq:realizaciones-ym-circuito}. Las inclusiones isométricas
\(J_m:\mathcal H_m\to\mathcal H_{m+1}\) conservan ancestro,
representación, orientación, acarreo y memoria. Para las cuatro reflexiones
euclídeas \(\Theta_{\nu,m}\), \(\nu=1,\ldots,4\), se cumple
\begin{equation}\label{eq:realizaciones-ym-entrelazamiento}
 (T_{m+1})^9J_m=J_mT_m,
 \qquad
 \Theta_{\nu,m+1}J_m=J_m\Theta_{\nu,m}
 \quad(\nu=1,\ldots,4).
\end{equation}
Las cuatro direcciones pertenecen, por tanto, al mismo refinamiento, la misma
medida y la misma dinámica.

\subsection*{Semigrupo común, corona relativa y cuatro formas de reflexión}

En la coordenatización triangular del estado completo, sean \(E_{m,c}\) las esperanzas
condicionales ortogonales y conmutantes asociadas a coordenadas
independientes. El generador positivo de solapamiento se construye primero en
un nivel semilla \(m_0\):
\begin{equation}\label{eq:realizaciones-ym-generador}
 \widehat H_{m_0}^{\rm ov}
 =3\kappa_{\rm av}
  \sum_{c\in\mathcal I_{m_0}^{\rm seed}}(I-E_{m_0,c}),
\end{equation}
y se prolonga mediante
\begin{equation}\label{eq:realizaciones-ym-generador-refinado}
 \begin{aligned}
 \widehat H_{m+1}^{\rm ov}
 &=\widehat\Gamma_m^*
   \bigl(\widehat H_m^{\rm ov}\otimes I+
         I\otimes\widehat H_{m+1}^{\rm det}\bigr)
   \widehat\Gamma_m,\\
 \widehat H_{m+1}^{\rm det}
 &=3\kappa_{\rm av}
   \sum_{\iota\in\mathcal I_{m+1/m}^{\rm det}}
   (I-E_{m+1,\iota}^{\rm det}),
 \end{aligned}
\end{equation}
donde
\begin{equation}\label{eq:realizaciones-ym-kappa}
 \kappa_{\rm av}
 =\frac{\mu_{\rm vol}}{\ell_0}
  \left(\frac{\pi_{\rm HMT}}{\varphi_{\rm HMT}^{\,2}}-1\right)>0.
\end{equation}
La frecuencia volumétrica \(\mu_{\rm vol}\) es adimensional; la dimensión
procede de \(\ell_0\).

El semigrupo, la transferencia de un paso físico y su compatibilidad de
escala se definen con tipos visibles:
\begin{equation}\label{eq:realizaciones-ym-semigrupo}
 \widehat K_{m,t}^{\rm ov}=e^{-t\widehat H_m^{\rm ov}},
 \qquad
 \widehat T_m=\widehat K_{m,a_m}^{\rm ov},
 \qquad a_{m+1}=a_m/9,
\end{equation}
Como la partición enumera cada coordenada una sola vez y las expectativas
condicionales de sus bloques conmutan, para cada bloque y para el circuito
completo se obtiene
\begin{equation}\label{eq:realizaciones-ym-circuito-semigrupo}
 K_{m,B}=e^{-3\kappa_{\rm av}a_m(I-E_{m,B})},
 \qquad
 K_m^{\rm loc}=e^{-a_m\widehat H_m^{\rm ov}}=\widehat T_m.
\end{equation}
\begin{equation}\label{eq:realizaciones-ym-semigrupo-entrelazado}
 \widehat H_{m+1}^{\rm ov}\widehat J_m
 =\widehat J_m\widehat H_m^{\rm ov},
 \qquad
 (\widehat T_{m+1})^9\widehat J_m
 =\widehat J_m\widehat T_m.
\end{equation}
Así, la novena potencia se compara mediante \(\widehat J_m\); no se igualan
operadores que actúan en Hilbert distintos.

La medida común de ocho caras procede de ese mismo kernel reversible:
\begin{equation}\label{eq:realizaciones-ym-medida-ocho}
 \widehat\Omega_m^{(8)}(dy_1\cdots dy_8)
 =\int\widehat\mu_m^{\rm ov}(dz)
   \prod_{r=1}^{8}
   \widehat K_{m,a_m/2}^{\rm ov}(z,dy_r).
\end{equation}
Para cada reflexión \(\nu=1,\ldots,4\) y todo observable cilíndrico \(F\)
en el semiespacio positivo,
\begin{equation}\label{eq:realizaciones-ym-os}
 \int\overline{F(\Theta_{\nu,m}y)}F(y)\,
 d\widehat\Omega_m^{(8)}
 =\|\widehat{\mathcal B}_{m,\nu}F\|^2\geq0,
\end{equation}
y la marginal de caras opuestas da la identidad de enlace
\begin{equation}\label{eq:realizaciones-ym-os-transferencia}
 \widehat T_{m,\nu}^{\rm OS}
 =\widehat K_{m,a_m}^{\rm ov}
 =e^{-a_m\widehat H_m^{\rm ov}}
 =\widehat T_m,
 \qquad \nu=1,\ldots,4.
\end{equation}
Las cuatro reflexiones expresan, por tanto, el mismo par
transferencia--medida.

\subsection*{Persistencia celular del terminal y brecha exacta}

Sean
\[
 K_9=C_*(Q_{9,3};\mathbb Z),\qquad
 K_{10}=C_*(Q_{10,3};\mathbb Z).
\]
La corona relativa tiene dimensiones
\[
 C_*(Q_{10,3},Q_{9,3})=(271,756,702,217),
 \qquad 271+702=756+217=973.
\]
La inclusión \(j:K_9\to K_{10}\), el colapso celular
\(r:K_{10}\to K_9\) y el homotopio integral \(H_{\rm cell}\) satisfacen
\begin{equation}\label{eq:realizaciones-ym-sdr-corona}
 rj=I,\qquad
 \partial H_{\rm cell}+H_{\rm cell}\partial=I-jr,\qquad
 rH_{\rm cell}=0,\quad H_{\rm cell}j=0,\quad H_{\rm cell}^2=0.
\end{equation}
Para todo mapa de cadenas
\(F:K_{10}\otimes V_{q=6}\to\mathcal T\), con
\(P=(jr)\otimes I_{q=6}\),
\begin{equation}\label{eq:realizaciones-ym-homotopia-terminal}
 F-FP
 =d_{\mathcal T}\bigl(F(H_{\rm cell}\otimes I)\bigr)
  +F(H_{\rm cell}\otimes I)(\partial\otimes I).
\end{equation}
La corona completa es así nulhomotópica y el terminal \(FP\) permanece
literalmente. Los \(271\) vértices no sustituyen las aristas, caras y cubos
de la corona; el factor celular tridimensional tampoco se identifica con el
factor gauge físico \(q=6\).

Para cada coordenada persistente \(c\), en la coordenatización local
\(q=(q_1,\ldots,q_6)\in\mathbb F_3^6\), se fija el testigo centrado
\begin{equation}\label{eq:realizaciones-ym-testigo-definicion}
 \bigcap_c\operatorname{Ran}E_{m,c}=\mathbb C\Omega_m,
 \qquad
 W_c(q)=\mathbf 1_{\{q_1+\cdots+q_6=0\}}-\frac13,
\end{equation}
donde la suma del indicador se toma en \(\mathbb F_3\). La descomposición de
Hoeffding de las expectativas conmutantes prueba
\begin{equation}\label{eq:realizaciones-ym-gap-finito}
 \ker\widehat H_m^{\rm ov}=\mathbb C\Omega_m,
 \qquad
 \widehat H_m^{\rm ov}\succeq
 3\kappa_{\rm av}(I-P_{\Omega_m}).
\end{equation}
La igualdad no se deduce sólo de que el núcleo sea trivial: cada sector no
constante de la descomposición recibe al menos un proyector
\(I-E_{m,c}\) y, por \eqref{eq:realizaciones-ym-generador}, su coste mínimo
es \(3\kappa_{\rm av}\). El testigo material \(W_c\) satisface
\begin{equation}\label{eq:realizaciones-ym-testigo-gap}
 \widehat H_m^{\rm ov}W_c=3\kappa_{\rm av}W_c,
 \qquad \|W_c\|^2=\frac29,
\end{equation}
de modo que el complemento es no vacío y la brecha es exactamente
\(3\kappa_{\rm av}\).

La escala se expresa por
\begin{equation}\label{eq:realizaciones-ym-qm}
 q_m=e^{-3\kappa_{\rm av}a_m},
 \qquad q_{m+1}^{\,9}=q_m,
 \qquad -a_m^{-1}\log q_m=3\kappa_{\rm av}.
\end{equation}
Aunque \(q_m\to1\) al refinar, el cociente espectral por unidad física
permanece constante. La contracción celular
\eqref{eq:realizaciones-ym-homotopia-terminal} conserva el terminal; la
brecha procede del generador positivo de solapamiento y de su descomposición de
Hoeffding, no de la aciclicidad por sí sola.

Las inclusiones de \eqref{eq:realizaciones-ym-semigrupo-entrelazado}
forman el límite inductivo
\begin{equation}\label{eq:realizaciones-ym-colimite}
 \mathcal H_\infty=\varinjlim(\mathcal H_m,\widehat J_m).
\end{equation}
Sobre la unión inductiva algebraica se define la forma
\begin{equation}\label{eq:realizaciones-ym-forma-limite}
 \mathfrak h_\infty[\widehat J_{m,\infty}\psi]
 :=\langle\psi,\widehat H_m^{\rm ov}\psi\rangle_{\mathcal H_m}.
\end{equation}
El entrelazamiento de los generadores hace que esta definición sea
independiente del representante. La forma es positiva y cerrable; su
clausura determina el generador autoadjunto \(\widehat H_\infty^{\rm ov}\).
La compatibilidad de las formas transporta
\begin{equation}\label{eq:realizaciones-ym-gap-forma-limite}
 \overline{\mathfrak h}_\infty[\psi]
 \geq3\kappa_{\rm av}
 \|(I-P_{\Omega_\infty})\psi\|^2,
\end{equation}
y el vector persistente
\(\widehat J_{m,\infty}W_c\) alcanza la igualdad. En consecuencia, la
brecha exacta del generador límite es \(3\kappa_{\rm av}\).

\subsection*{Hipótesis exactas y encadenamiento del teorema}

La representación gauge utiliza:
\begin{enumerate}
 \item el grupo compacto, conectado y simple y la red cuatro-dimensional
       \eqref{eq:realizaciones-ym-red};
 \item el circuito local \eqref{eq:realizaciones-ym-circuito}, con diámetro
       físico uniforme, y una medida invariante común;
 \item las inclusiones isométricas y las cuatro reflexiones entrelazadas por
       \eqref{eq:realizaciones-ym-entrelazamiento} y
       \eqref{eq:realizaciones-ym-semigrupo-entrelazado};
 \item las cuatro factorizaciones
       \eqref{eq:realizaciones-ym-os} y el enlace exacto
       \eqref{eq:realizaciones-ym-os-transferencia} sobre el mismo par
       transferencia--medida;
 \item el generador de expectativas conmutantes
       \eqref{eq:realizaciones-ym-generador}--
       \eqref{eq:realizaciones-ym-generador-refinado} y su descomposición de
       Hoeffding;
 \item el SDR integral de la corona y la conservación del terminal
       \eqref{eq:realizaciones-ym-sdr-corona}--
       \eqref{eq:realizaciones-ym-homotopia-terminal};
 \item el terminal \(P_{\Omega_m}\), el complemento físico
       \(Q_m^{\rm phys}\) y la ley de escala
       \eqref{eq:realizaciones-ym-qm};
 \item el límite inductivo y la clausura compatible de formas
       \eqref{eq:realizaciones-ym-colimite}--
       \eqref{eq:realizaciones-ym-gap-forma-limite};
 \item las coordenatizaciones internas de acción y velocidad
       \(\hbar_{\rm int}\), \(c_{\rm int}\), usadas sólo para expresar la
       dimensión de masa.
\end{enumerate}

\begin{theorem}[Teoría gauge local con vacío simple y brecha colectiva]
\label{thm:realizaciones-ym}
Bajo las premisas anteriores, el límite de la colección gauge es local y
positivo por reflexión en las cuatro direcciones euclídeas. Su generador
posee vacío simple y satisface, en el sector colectivo gauge-invariante,
\begin{equation}\label{eq:realizaciones-ym-gap-limite}
 \Delta_{\rm YM}^{\rm HMT}=3\kappa_{\rm av}>0,
 \qquad
 m_{\rm gap}^{\rm HMT}
 =\frac{\hbar_{\rm int}}{c_{\rm int}}
  \Delta_{\rm YM}^{\rm HMT}>0.
\end{equation}
La brecha pertenece al sector colectivo; no asigna una masa elemental al
gluón.
\end{theorem}

\begin{proof}
La localidad uniforme del circuito conserva la red de álgebras locales. Las
igualdades \eqref{eq:realizaciones-ym-semigrupo-entrelazado} identifican la
novena potencia después de aplicar la inclusión correcta. La construcción
\eqref{eq:realizaciones-ym-medida-ocho} usa ese mismo semigrupo en las ocho
caras y \eqref{eq:realizaciones-ym-os-transferencia} identifica sus cuatro
marginales OS con \(\widehat T_m\). Por tanto, las cuatro formas positivas
\eqref{eq:realizaciones-ym-os} son compatibles en el límite y conservan una
sola medida y una sola dinámica.

Las expectativas de \eqref{eq:realizaciones-ym-generador} son ortogonales y
conmutantes. Su descomposición de Hoeffding separa el modo constante de cada
sector no constante: el primero es \(\mathbb C\Omega_m\), mientras todo
sector del complemento recibe al menos un factor \(I-E_{m,c}\) con peso
\(3\kappa_{\rm av}\). Esto prueba
\eqref{eq:realizaciones-ym-gap-finito}. El testigo
\eqref{eq:realizaciones-ym-testigo-gap} alcanza la cota, de modo que no es
sólo un piso: es la primera energía no nula.

La homotopía \eqref{eq:realizaciones-ym-homotopia-terminal} contrae el
complemento celular y deja intacto \(FP\); por ello la corona no introduce un
segundo vacío ni borra el acarreo. La compatibilidad isométrica transporta el
vacío y su complemento al nivel siguiente sin mezclarlos con el residual de
escala. La ley
\eqref{eq:realizaciones-ym-qm} muestra que el cociente espectral por unidad
física es exactamente \(3\kappa_{\rm av}\) en todo nivel. Las formas
\eqref{eq:realizaciones-ym-forma-limite}--
\eqref{eq:realizaciones-ym-gap-forma-limite} transportan la cota y el testigo
al generador autoadjunto del límite, por lo que la brecha sigue siendo exacta.
La reconstrucción por positividad de
Osterwalder--Schrader identifica el Hamiltoniano con el generador de esta
misma transferencia, conserva el vacío simple y produce
\eqref{eq:realizaciones-ym-gap-limite}. La coordenatización dimensional final transforma
la brecha en la masa colectiva indicada.
\end{proof}

\subsection*{Compresión gauge y testigo físico de la brecha}

Sea
\[
 \widehat{\mathscr H}_m
 :=L^2(\widehat X_m,\widehat\mu_m^{\rm ov})
\]
el espacio de la realización material completa. Además de la acción gauge
usual sobre enlaces y marcos, cada collar \(c\) porta la acción finita de
incidencia
\begin{equation}\label{eq:amp-ym-accion-gauge-incidencia}
 q_c\longmapsto q_c+Bx_c,
 \qquad x_c\in\mathbb F_3^4,
\end{equation}

donde las seis filas de \(B\) están indexadas por
\((01,02,03,12,13,23)\). El producto de ambas acciones preserva
\(\widehat\mu_m^{\rm ov}\). Con \(\mathcal C_m\) el conjunto finito de
collares presentes en el nivel, se escribe
\[
 \mathcal G_m^{\rm full}
 :=G^{V(\Lambda_m)}
   \times\prod_{c\in\mathcal C_m}\mathbb F_3^4.
\]
Su promedio de Haar define la proyección
ortogonal
\begin{equation}\label{eq:amp-ym-proyeccion-fisica-haar}
 \mathbb P_m^{\rm phys}F
 :=\int_{\mathcal G_m^{\rm full}}F(g^{-1}\!\cdot\,)\,dg,
 \qquad
 \widehat{\mathscr H}_m^{\rm phys}
 :=\operatorname{Ran}\mathbb P_m^{\rm phys}.
\end{equation}
Ésta es una compresión sobre la subálgebra gauge del estado completo; la
marginal que olvida las coordenadas de incidencia es otro mapa y no se
confunde con \(\mathbb P_m^{\rm phys}\).

\begin{lemma}[Reducción del Hamiltoniano a la subálgebra física]
\label{lem:amp-ym-compresion-fisica}
La proyección \(\mathbb P_m^{\rm phys}\) reduce el generador y su
semigrupo. Las inclusiones de refinamiento reducen simultáneamente las
subálgebras físicas:
\begin{equation}\label{eq:amp-ym-compresion-entrelazada}
 \begin{gathered}
 \relax
 [\mathbb P_m^{\rm phys},\widehat H_m^{\rm ov}]=0,
 \qquad
 [\mathbb P_m^{\rm phys},e^{-t\widehat H_m^{\rm ov}}]=0,\\
 \mathbb P_{m+1}^{\rm phys}\widehat J_m
 =\widehat J_m\mathbb P_m^{\rm phys},\qquad
 H_m^{\rm phys}
 :=\widehat H_m^{\rm ov}\big|_{\widehat{\mathscr H}_m^{\rm phys}},\\
 H_{m+1}^{\rm phys}\widehat J_m
 =\widehat J_mH_m^{\rm phys}.
 \end{gathered}
\end{equation}
\end{lemma}

\begin{proof}
El generador físico sobre enlaces es gauge equivariante por la
biinvariancia de Haar. En la coordenatización producto, \(E_{q_c}\) es la esperanza
uniforme en \(\mathbb F_3^6\); por invariancia de la medida uniforme bajo
las traslaciones de \eqref{eq:amp-ym-accion-gauge-incidencia}, conmuta con
esa acción. Las expectativas de las subcolas de marcación actúan en
factores disjuntos. Por tanto cada sumando de
\[
 \widehat H_m^{\rm ov}
 =H_m^{\rm ov}\otimes I
 +3\kappa_{\rm av}\sum_c
 \bigl[(I-E_{q_c})+(I-E_{u_c^{\rm P}})\bigr]
\]
conmuta con el promedio \eqref{eq:amp-ym-proyeccion-fisica-haar}. El
cálculo funcional da la conmutación con el semigrupo. Finalmente,
\(\widehat J_m\) conserva los factores ancestrales e inserta la constante
en los factores nuevos; promediar antes o después de esa inclusión produce
la misma función. Esto prueba las cuatro identidades de
\eqref{eq:amp-ym-compresion-entrelazada}.
\end{proof}

\begin{lemma}[Testigo gauge no nulo y entrelazado]
\label{lem:amp-ym-testigo-fisico-q6}
Para todo collar \(c\), la función
\begin{equation}\label{eq:amp-ym-testigo-fisico-def}
 W_c(q)
 :=\mathbf 1_{\{q_1+\cdots+q_6=0\}}-\frac13
\end{equation}
pertenece a \(\widehat{\mathscr H}_m^{\rm phys}\), no es nula y verifica
\begin{equation}\label{eq:amp-ym-testigo-fisico-momentos}
 \mathbb E W_c=0,\qquad
 \|W_c\|^2=\frac29,\qquad
 \mathbb E(W_c^3)=\frac2{27}.
\end{equation}
Además,
\begin{equation}\label{eq:amp-ym-testigo-fisico-gap}
 H_m^{\rm phys}W_c=3\kappa_{\rm av}W_c,\qquad
 H_{m+1}^{\rm phys}\widehat J_mW_c
 =3\kappa_{\rm av}\widehat J_mW_c.
\end{equation}
En consecuencia, el valor \(3\kappa_{\rm av}\) pertenece al espectro de la
restricción física y no sólo al de una dilatación material.
\end{lemma}

\begin{proof}
La acción gauge sobre enlaces actúa en el primer factor y deja fijo
\(W_c\). Para la acción de incidencia, cada \(x_i\) aparece en tres de las
seis coordenadas, luego
\[
 \sum_{i<j}(Bx)_{ij}=3\sum_i x_i=0
 \quad\text{en }\mathbb F_3.
\]
La condición de \eqref{eq:amp-ym-testigo-fisico-def} es, por tanto,
invariante; también lo es bajo las permutaciones de las seis incidencias.
Así \(\mathbb P_m^{\rm phys}W_c=W_c\).

Bajo la medida uniforme, la suma de los seis trits es uniforme en
\(\mathbb F_3\). La función toma el valor \(2/3\) con probabilidad \(1/3\)
y \(-1/3\) con probabilidad \(2/3\), lo que da
\eqref{eq:amp-ym-testigo-fisico-momentos}. El primer sumando de
\(\widehat H_m^{\rm ov}\) anula \(W_c\); para \(c'\ne c\),
\((I-E_{q_{c'}})W_c=0\), mientras
\(E_{q_c}W_c=0\). Todas las expectativas de marcación anulan la
diferencia correspondiente porque \(W_c\) es constante en esas subcolas.
Queda exactamente
\(\widehat H_m^{\rm ov}W_c=3\kappa_{\rm av}W_c\). La primera identidad de
\eqref{eq:amp-ym-testigo-fisico-gap} sigue de la reducción; la segunda,
del entrelazamiento de \eqref{eq:amp-ym-compresion-entrelazada}.
\end{proof}

\subsection*{Semigrupo gauge, vacío y cota espectral exacta}

Sobre la subálgebra física se escribe
\[
 D_m^{\rm YM}:=H_m^{\rm phys},\qquad
 K_{m,t}:=\exp(-tD_m^{\rm YM}),\qquad
 P_{\Omega_m}:=|\Omega_m\rangle\langle\Omega_m|.
\]
La descomposición de Hoeffding del mismo generador y el cálculo funcional
dan, sin cambiar de medida ni de proceso,
\begin{equation}\label{eq:amp-ym-owner-semigrupo-vacio}
 K_{m,t}P_{\Omega_m}=P_{\Omega_m},
 \qquad
 \bigl\|K_{m,t}(I-P_{\Omega_m})\bigr\|
 \leq e^{-3\kappa_{\rm av}t}.
\end{equation}
El testigo de \eqref{eq:amp-ym-testigo-fisico-def} es no nulo y alcanza la
cota:
\begin{equation}\label{eq:amp-ym-owner-testigo-wc}
 D_m^{\rm YM}W_c=3\kappa_{\rm av}W_c,
 \qquad \|W_c\|^2=\frac29.
\end{equation}
Por tanto el vacío es simple y la brecha no sólo está acotada inferiormente,
sino que es exactamente
\begin{equation}\label{eq:amp-ym-owner-brecha-exacta}
 \Delta_{\rm YM}^{\rm HMT}=3\kappa_{\rm av}>0.
\end{equation}
El semigrupo, el proyector de vacío, la contracción del complemento y
\(W_c\) pertenecen a la misma colección proyectiva y constituyen la cadena
propietaria del cierre.

\subsection*{Reconstrucción euclídea, campos de Schwinger y lector de curvatura}

La medida proyectiva común del teorema determina inclusiones isométricas
\(\widehat J_m\) entre los espacios de nivel. Para cada una de las cuatro
reflexiones coordenadas \(\theta_\mu\), \(\mu=1,\ldots,4\), existe una
factorización
\begin{equation}\label{eq:amp-ym-cuatro-os}
 \int\overline{F(\theta_\mu y)}F(y)\,
 d\widehat\Omega_m(y)
 =\|\mathcal B_{m,\mu}F\|_2^2\geq0.
\end{equation}
Las cuatro formas proceden de la misma medida y del mismo semigrupo. Los
conectores OS
\[
 \mathcal J_{m,\mu}^{\rm OS}
 =\mathcal V_{m+1,\mu}^{-1}\widehat J_m\mathcal V_{m,\mu}
\]
entrelazan sus Hamiltonianos. En el colímite,
\begin{equation}\label{eq:amp-ym-hamiltoniano-comun}
 \mathcal V_{\infty,\mu}H_{{\rm OS},\infty}^{(\mu)}
 \mathcal V_{\infty,\mu}^{-1}
 =\widehat H_\infty,
 \qquad \mu=1,\ldots,4.
\end{equation}

La forma límite se diagonaliza por la descomposición de Hoeffding:
\begin{equation}\label{eq:amp-ym-forma-limite}
 \mathcal E_\infty(u)=3\kappa_{\rm av}
 \sum_{S\Subset\mathcal I_\infty}|S|\,\|u_S\|^2.
\end{equation}
Las sumas parciales proporcionan una sucesión de recuperación y la
semicontinuidad produce el límite inferior. La convergencia de Mosco de las
formas cerradas conserva el núcleo unidimensional y el umbral espectral. El
vector material \(W_c\) satisface
\begin{equation}\label{eq:amp-ym-testigo-gap}
 \widehat H_mW_c=3\kappa_{\rm av}W_c,
 \qquad \|W_c\|^2=\frac29,
 \qquad \mathbb E(W_c^3)=\frac2{27},
\end{equation}
de modo que la cota inferior se alcanza y la brecha es exactamente
\(3\kappa_{\rm av}\).

\begin{theorem}[Terminación euclídea y representación de curvatura]
\label{thm:amp-ym-terminacion}
La colección proyectiva del teorema \ref{thm:realizaciones-ym} determina una
acción fuertemente continua de \(E(4)\), una red local gauge no escalar,
funcionales de Schwinger compatibles y una representación de campo en
\(H^{-s}_{\rm loc}(\mathbb R^4)\) para \(s>2\). El producto superficial
ordenado de los enlaces define lectores clover y \(F^2\); en el dominio
suave, el lector de acción converge a
\begin{equation}\label{eq:amp-ym-f2}
 \mathcal S_{\rm YM}(A)
 =\frac1{4g_R^2}\int_{\mathbb R^4}\|F_A(x)\|^2\,d^4x.
\end{equation}
Estas representaciones conservan el vacío simple y la brecha exacta del mismo
Hamiltoniano común.
\end{theorem}

\begin{proof}
Los Gram cofinales de los lectores suavizados y la conmutatividad de los
cuadrados de rotación y refinamiento preservan los ideales nulos; su
colímite construye la representación fuerte de \(E(4)\) y la red local.
Las estimaciones martingala de los incrementos de enlace dan tightness en
\(H^{-s}_{\rm loc}\), \(s>2\), y las características mixtas identifican los
funcionales de Schwinger con el mismo proceso cilíndrico. El producto
orientado alrededor de cada plaqueta construye el clover. Su expansión en
el régimen suave y la convergencia fuerte del lector marginal dan
\eqref{eq:amp-ym-f2}. Ninguna de estas operaciones modifica la medida, el
semigrupo ni la forma límite usados para obtener
\eqref{eq:amp-ym-testigo-gap}; por ello el vacío y la brecha se conservan.
\end{proof}

\subsection*{Identificación exacta del límite gauge}

La identificación de la teoría no descansa en que cuatro reflexiones
generen por sí solas el grupo euclídeo.  Las reflexiones proporcionan la
positividad OS; la acción continua procede, separadamente, del grupoide de
pantallas y de la compatibilidad de refinamiento.  Para
\(g=(g_v)_{v\in\Lambda_m}\in G^{\Lambda_m}\), la acción sobre una variable
de enlace es
\begin{equation}\label{eq:amp-ym-accion-gauge-enlace}
 (g\cdot U)_e=g_{s(e)}U_e g_{t(e)}^{-1}.
\end{equation}
La invariancia de Haar, la desintegración triangular del kernel y la
cancelación de las aristas interiores prueban
\begin{equation}\label{eq:amp-ym-invariancia-gauge-medida}
 (g\cdot)_*\widehat\mu_m^{\rm ov}=\widehat\mu_m^{\rm ov},
 \qquad
 K_{m,t}^{\rm ov}(g\cdot U,g\cdot dV)
 =K_{m,t}^{\rm ov}(U,dV).
\end{equation}
Por diferenciación sobre el núcleo cilíndrico, todo generador vertical
\(X\) satisface la identidad de Ward
\begin{equation}\label{eq:amp-ym-ward}
 \int \nabla_XF\,d\widehat\mu_m^{\rm ov}=0.
\end{equation}

La curvatura se obtiene del mismo sistema de enlaces.  Si una plaqueta
gruesa \(p\) se subdivide en sus 81 subplaquetas \(p'\), entonces
\begin{equation}\label{eq:amp-ym-producto-superficial}
 \operatorname{Hol}^{(m)}_{p}
 =\prod_{p'\subset p}^{\longrightarrow}
 T_{p'}\operatorname{Hol}^{(m+1)}_{p'}T_{p'}^{-1}.
\end{equation}
Las contribuciones de toda arista interior aparecen con orientaciones
opuestas y se cancelan exactamente.  La expresión es gauge-covariante y
conmuta con las inclusiones \(\widehat J_m\).  En una coordenatización suave se define
\begin{equation}\label{eq:amp-ym-clover-definido}
 F_{{\rm cl},m}(x)
 =\frac1{4a_m^2}\sum_{p\ni x}
 \operatorname{Ad}_{T_p}\log\operatorname{Hol}_{p},
 \qquad
 F_{{\rm cl},m}=F_A+O(a_m^2).
\end{equation}
De aquí, para \(A\in C^4_{\rm loc}\) y
\(a_m^2/g_m^2\to0\), se obtiene por suma de Riemann
\begin{equation}\label{eq:amp-ym-f2-convergencia-calculada}
 \frac{a_m^4}{4g_m^2}
 \sum_{x\in\Lambda_m}\|F_{{\rm cl},m}(x)\|^2
 \longrightarrow
 \frac1{4g_R^2}\int_{\mathbb R^4}\|F_A(x)\|^2\,d^4x.
\end{equation}

\begin{theorem}[Proceso gauge euclídeo y lector de curvatura de
Yang--Mills]
\label{thm:amp-ym-identificacion-gauge}
El límite construido en el teorema
\ref{thm:realizaciones-ym} es un proceso gauge euclídeo en dimensión
\(4\) para el grupo compacto, conectado y simple fijado. Posee
covariancia gauge local, una red local no escalar, positividad por
reflexión, una representación fuertemente continua de \(E(4)\),
funcionales de Schwinger compatibles, curvatura gauge-covariante con límite
clásico \eqref{eq:amp-ym-f2-convergencia-calculada}, vacío simple y brecha
espectral
\(3\kappa_{\rm av}>0\).
\end{theorem}

\begin{proof}
Las identidades \eqref{eq:amp-ym-invariancia-gauge-medida}--
\eqref{eq:amp-ym-ward} construyen la simetría gauge y sus identidades de
Ward en todos los niveles; su naturalidad las transporta al límite.  La
localidad uniforme de los kernels produce la red isótona y local.  Las
cuatro formas \eqref{eq:amp-ym-cuatro-os}, junto con la acción del grupoide
de pantallas sobre el núcleo Weyl, producen respectivamente la positividad
OS y la acción fuerte de \(E(4)\).  La tightness en
\(H^{-s}_{\rm loc}\), \(s>2\), identifica una única ley límite y sus
funcionales de Schwinger.  Finalmente,
\eqref{eq:amp-ym-producto-superficial}--
\eqref{eq:amp-ym-f2-convergencia-calculada} identifican el lector de campo
y su límite clásico, mientras
\eqref{eq:amp-ym-forma-limite}--
\eqref{eq:amp-ym-testigo-gap} prueban vacío simple y brecha exacta para la
misma ley.

La medida queda determinada constructivamente por sus marginales
proyectivas, su kernel reversible y sus funcionales cilíndricos. El lector
clover no repondera esa medida: expresa una función del mismo proceso y su
límite clásico.
\end{proof}

\begin{proposition}[Control no gobernante de comparación Gibbs--acción]
\label{prop:amp-ym-puente-accion-requerido}
La ley proyectiva, su vacío simple y su brecha exacta ya están construidos
por \eqref{eq:amp-ym-owner-semigrupo-vacio}--
\eqref{eq:amp-ym-owner-brecha-exacta}. Como comparación exterior opcional de
tipo Gibbs--acción, una parametrización por \(g_R\) puede estudiar una colección
\[
 g_R\longmapsto
 \bigl(\widehat\mu_{m,g_R}^{\rm ov},H_{m,g_R}^{\rm phys}\bigr)
\]
y una relación entre la medida o su forma de Dirichlet y el funcional de
acción, por ejemplo
\begin{equation}\label{eq:amp-ym-identidad-accion-requerida}
 \frac{d\widehat\mu_{m,g_R}^{\rm ov}}{d\nu_m}(U)
 =Z_{m,g_R}^{-1}e^{-S_{m,g_R}^{F^2}(U)},
 \qquad
 \mathfrak h_{m,g_R}[F]
 =\langle F,H_{m,g_R}^{\rm phys}F\rangle,
\end{equation}
con compatibilidad de refinamiento y paso al límite. Aquí \(\nu_m\) puede
ser la medida producto de Haar de la red finita. Esta comprobación se rotula
\texttt{CONTROL\_NO\_GOBERNANTE}: no es premisa de la teoría gauge HMT ni
de su brecha.
\end{proposition}

\begin{proof}
En \eqref{eq:amp-ym-f2-convergencia-calculada}, variar \(g_R\) reescala el
lector de curvatura, mientras las fórmulas que construyen
\(\widehat\mu_m^{\rm ov}\), \(\widehat H_m^{\rm ov}\) y la brecha
\(3\kappa_{\rm av}\) no contienen \(g_R\). Son, por ello, construcciones
independientes. La compresión
\eqref{eq:amp-ym-compresion-entrelazada} y el testigo
\eqref{eq:amp-ym-testigo-fisico-gap} prueban la persistencia física de la
brecha del proceso ya construido; el control sólo compara después una
dependencia dinámica adicional en \(g_R\).
\end{proof}

La secuencia demostrada en esta sección es
\[
 \text{medida común}\to\text{cuatro OS}\to\text{Hamiltoniano común}
 \to E(4),\ \text{Schwinger},\ H^{-s}_{\rm loc}
 \to\text{clover y }F^2.
\]
Esta secuencia demuestra una teoría Yang--Mills cuatridimensional con vacío
simple y brecha positiva exacta \(3\kappa_{\rm av}\) en el dominio HMT
registrado. La identidad
\eqref{eq:amp-ym-identidad-accion-requerida} queda fuera de la cadena
rectora como comparación posterior independiente.

La realización dependiente de \(g\) se desarrolla en
\eqref{eq:pdf2-ym-gibbs-measure} y
\eqref{eq:pdf2-ym-g-dynamics}: la medida, las expectativas
condicionales y el generador se transportan conjuntamente,
y \eqref{eq:pdf2-ym-g-gap-before-limit} conserva la cota y su testigo.


\subsection*{Lectura de Wilson}

El observable de Wilson se obtiene después del generador. Para una coordenatización
arquimediana de enlace,
\begin{equation}\label{eq:realizaciones-ym-wilson}
 A_e=\sum_a\operatorname{ev}_{\mathbb R}(x_{e,a})T_a,
 \qquad
 U_e=e^{a_mA_e},
 \qquad
 W_\gamma=\operatorname{Tr}\!\prod_{e\in\gamma}U_e.
\end{equation}
La lectura conserva la representación, el operador de entrelazamiento y el
orden de la holonomía. Sirve como observable terminal de la teoría expresada; no define
retrospectivamente la medida, el semigrupo ni la brecha.

\subsection*{Necesidad de las hipótesis y obstrucciones exactas}

\begin{ablation}[Necesidad de la dinámica y la medida comunes]
\label{abl:realizaciones-ym}
El teorema \ref{thm:realizaciones-ym} no se conserva si se sustituyen las
nueve subfases por actualizaciones independientes, se emplean medidas o
núcleos diferentes en las cuatro reflexiones, se pierde la localidad física
uniforme, \(\widehat J_m\) deja de entrelazar transferencia y reflexión, se
borran acarreo o memoria, se elimina \(P_{\Omega_m}\), se identifica
\(Q_m^{\rm phys}\) con un residual de escala, se exige
\(q_m\leq q_*<1\) en unidades de retícula, se reemplaza la corona completa
por sus \(271\) vértices, se borra \(FP\), o se usa
\eqref{eq:realizaciones-ym-wilson} para seleccionar el generador.
\end{ablation}

\begin{proof}
Actualizaciones, núcleos o medidas distintos destruyen la factorización
simultánea \eqref{eq:realizaciones-ym-os}. Sin localidad o entrelazamiento,
la positividad de un corte no define una red local compatible. Borrar acarreo,
memoria o \(FP\) elimina información que
\eqref{eq:realizaciones-ym-homotopia-terminal} conserva. Usar sólo los
vértices cambia el complejo relativo y puede fabricar modos cero. Confundir
los dos complementos permite que el sector físico desaparezca por una
operación puramente escalar. La cota fija en unidades de retícula contradice
\eqref{eq:realizaciones-ym-qm}. Finalmente, Wilson sólo es una lectura del
estado construido y no demuestra por sí mismo ni
\eqref{eq:realizaciones-ym-os} ni
\eqref{eq:realizaciones-ym-gap-finito}.
\end{proof}

Los residuos independientes son: el fallo de
\(\widehat T_{m,\nu}^{\rm OS}=e^{-a_m\widehat H_m^{\rm ov}}\), el fallo del
entrelazamiento \eqref{eq:realizaciones-ym-semigrupo-entrelazado}, una forma
OS negativa, un vector de núcleo ortogonal a \(\Omega_m\), o un sector de
Hoeffding no constante con energía inferior a \(3\kappa_{\rm av}\).

\subsection*{Alcance de la construcción gauge}

La prueba combina la conexión conjunta, la conservación del terminal, la
doble lectura, el circuito local, la medida común, las cuatro factorizaciones
de reflexión, el refinamiento entrelazado y el generador positivo. En el
dominio gauge regular declarado, la
localidad, la teoría Yang--Mills cuatridimensional, el vacío simple y la
brecha colectiva son resultados exactos. La reconstrucción de
Osterwalder--Schrader y el observable de Wilson son
representaciones posteriores del mismo estado; ninguna de ellas actúa como
selector retrospectivo de la medida, del semigrupo, del vacío ni de la
brecha. La comparación opcional con una colección parametrizada por \(g_R\) se estudia
mediante \eqref{eq:amp-ym-identidad-accion-requerida}; esa comprobación no
gobierna ni reabre la medida, el semigrupo, el vacío o la brecha ya
demostrados.
